\begin{tomtat}
\begin{itemize}
  \item Ba bài báo của chương đứng ở ba chỗ khác nhau so với chữ trong nhan đề,
và gộp chúng lại là đọc sai cả ba: ở bài thứ nhất lời giải thích là nạn nhân, ở
bài thứ hai nó là vũ khí nhắm vào một bộ phân loại ảnh, còn ở bài thứ ba không
có bên tấn công nào và chỗ hỏng được sửa là một chỗ hỏng thống kê. Chữ
  \enquote{faithful} không xuất hiện lần nào trong cả ba bài, và trong bốn họ
chỉ số của chương~\ref{ch:do-do-trung-thuc} chỉ mỗi insertion có mặt, ở bài thứ
ba, nơi nó là dụng cụ dùng để chấm điểm phép sửa.
  \item \tn{Phép dựng giàn}{scaffolding} của phương trình~\ref{eq:ch14-phep-dung-gian} công bố một
mô hình trả về dự đoán thiên lệch trên dữ liệu thật và một mô hình vô hại trên
những điểm đã nhiễu, rồi để một
  \tn{bộ phát hiện ngoài phân phối}{out-of-distribution detector} chọn nhánh. Nó
không nhắm vào một giá trị tham số nào của LIME hay SHAP mà nhắm vào điều đúng
với mọi giá trị, rằng những điểm được hỏi không phải những điểm mô hình sẽ chạy
trên, nên đổi siêu tham số không gỡ được. Chỗ SHAP cầm cự khá hơn khi có hai
đặc trưng đánh lạc hướng đến từ độ chính xác cục bộ của định
nghĩa~\ref{def:ch04-ba-tinh-chat}, tức từ một tiên đề chứ không từ một biện
pháp phòng thủ.
  \item Thiết lập ấy dựng ra một ground truth cho attribution đặc trưng, và dựng
bằng đường thứ ba: không phải qua thiết kế nhiệm vụ như
chương~\ref{ch:chi-so-khong-do-do-trung-thuc}, cũng không phải qua thiết kế dữ
liệu, mà qua thiết kế chính mô hình, vì bên tấn công buộc phải biết mô hình của
mình dựa vào cái gì. Cái giá là tính đại diện: kết luận rút ra là kết luận về
một mô hình được dựng riêng để đánh bại phương pháp. Đọc theo
bảng~\ref{tab:ch01-bon-quan-he} thì phép tấn công không phá tính trung thực
dưới một tập can thiệp cố định, nó tách hai tập can thiệp ra rồi đẩy chúng đi
hai đường.
  \item Sáu chỗ chương này không nhận lập luận của nguồn. Bài thứ nhất phát biểu
ngưỡng sụp của LIME bằng ba giá trị khác nhau rồi báo cáo phương pháp của mình
suy giảm ở một giá trị nằm giữa ba giá trị ấy, và tự ghi trong phần công việc
tương lai rằng phép so ấy không công bằng; nó còn dựng phép ghép trên một câu
sai về khoảng giá trị của SHAP, và phép nhân mà câu ấy biện hộ làm mất dấu của
điểm số. Ở bài thứ hai, phương trình~\ref{eq:ch14-tan-cong-shap} chỉ làm một
trong hai chiều dịch chuyển mà lời văn quanh nó mô tả, vì nhánh được chọn theo
giá trị điểm ảnh chứ không theo dấu của điểm số. Ở bài thứ ba, định lý về tính
vắng mặt chứng minh một phát biểu khác với tiên đề nó mang tên, và câu nói rằng
tính nhất quán cũng đúng cho điểm số đã chuẩn hóa thì sai, vì hệ số chuẩn hóa
phụ thuộc vào mô hình còn tính nhất quán so hai mô hình.
  \item Cột nhãn RMSE trong bảng kết quả tổng hợp của bài thứ ba là tổng bình
phương sai lệch chứ không phải căn của trung bình bình phương: cách đọc ấy dựng
lại mười một trong mười hai ô, hai ô trong đó đúng trọn bốn chữ số thập phân.
Hai cách đọc đồng biến với nhau nên mọi thứ hạng giữ nguyên và kết luận của bài
đứng vững; cái sai là cái nhãn.
  \item Cả ba phương pháp trong chương đều là phương pháp hậu nghiệm, và cả ba
bài đều xoay quanh cùng một chỗ: điểm số chỉ có được bằng cách hỏi mô hình về
những đầu vào đã bị thay hoặc bị che, tức những đầu vào không phải dữ liệu
thật. Yêu cầu mà chương để lại cho một dụng cụ tương lai là yêu cầu công bố tập
can thiệp, và nó khớp với yêu cầu mà mục~\ref{sec:ch13-tran-noi-gi} rút ra từ
một cái trần lý thuyết.
\end{itemize}
\end{tomtat}

\cauhoi
\begin{cauhoids}
  \item Mục~\ref{sec:ch14-ground-truth-ke-tan-cong} lập luận rằng lời giải thích
mà LIME trả về cho mô hình $e$ không sai dưới tập can thiệp của chính LIME. Hãy
phát biểu lại định nghĩa~\ref{def:ch01-do-trung-thuc} với tập can thiệp viết
tường minh thành một đối số, rồi cho biết dưới cách phát biểu ấy thì phép dựng
giàn còn là một phản ví dụ về độ trung thực hay đã thành một phản ví dụ về một
thứ khác, và thứ khác ấy là gì.
  \item Bài của mục~\ref{sec:ch14-mot-phong-thu} ghép LIME với SHAP bằng phép
nhân. Hãy lấy một đặc trưng mà LIME cho điểm $-0.6$ và SHAP cho điểm $-0.5$,
một đặc trưng khác mà cả hai cho $0.6$ và $0.5$, rồi tính điểm số BASIC SHLIME
cho cả hai. Sau đó đề xuất một cách ghép khác giữ được dấu, và cho biết cách ấy
còn giữ được tính chất nào trong ba tính chất của định
nghĩa~\ref{def:ch04-ba-tinh-chat} hay không.
  \item Phương trình~\ref{eq:ch14-tan-cong-shap} chọn nhánh theo giá trị điểm
ảnh chứ không theo dấu của $\phi$. Hãy viết lại nó sao cho nó làm đúng cả hai
chiều dịch chuyển mà lời văn của bài mô tả, rồi cho biết phiên bản viết lại có
cần thêm giả thiết nào về quan hệ giữa giá trị điểm ảnh và dấu của $\phi$ hay
không.
  \item Mục~\ref{sec:ch14-shap-lam-vu-khi} đọc tỉ lệ phân loại sai như một phép
đo về bộ điểm số đã dẫn đường cho phép tấn công. Hãy chỉ ra phép đo ấy giống họ
xóa dần của mục~\ref{sec:ch08-ho-xoa-dan} ở những chỗ nào và khác ở những chỗ
nào, rồi cho biết phản đối về đầu vào ngoài phân phối ở
mục~\ref{sec:ch08-huan-luyen-lai} có áp được vào nó hay không, và vì sao.
  \item Giá trị tham chiếu trong nửa thí nghiệm tổng hợp của
mục~\ref{sec:ch14-sua-bang-nhan-qua} được lấy bằng cách chạy SHAP trên một tập
đặc trưng đã rút gọn. Hãy nêu một phương pháp attribution giả định sẽ đạt điểm
hoàn hảo theo cách chấm ấy mà vẫn không trung thực theo định
nghĩa~\ref{def:ch01-do-trung-thuc}, rồi cho biết phải đổi cách dựng giá trị
tham chiếu thế nào để loại nó.
  \item Sách bác câu nói rằng tính nhất quán vẫn đúng sau bước chuẩn hóa. Hãy
dựng một phản ví dụ bằng số: hai mô hình $f$ và $f'$, ba đặc trưng, các điểm số
trước chuẩn hóa và hai giá trị $f(x)$, $f'(x)$ sao cho một đặc trưng có điểm số
trước chuẩn hóa tăng mà điểm số sau chuẩn hóa giảm.
  \item Ba đường dựng ground truth mà mục~\ref{sec:ch14-ground-truth-ke-tan-cong}
liệt ra ép ba thứ khác nhau: thiết kế nhiệm vụ, thiết kế dữ liệu, thiết kế mô
hình. Hãy xếp ba đường theo mức độ mà kết luận rút ra từ chúng nói được về
những mô hình ngoài đời, rồi cho biết chương nào trong sách cung cấp bằng chứng
cho chỗ xếp của mỗi đường.
  \item Đọc phần \enquote{Detecting OOD Samples} trong \enquote{Fooling LIME and
SHAP: Adversarial Attacks on Post hoc Explanation Methods} của Slack và cộng
sự~\autocite{mfooling}, nơi mô tả cách dựng bộ phát hiện ngoài phân phối. Hãy
cho biết bên kiểm tra có thể làm gì với chính phân phối nhiễu của mình để phép
tấn công khó thành hơn, và vì sao cách ấy không phải một lời giải chung.
\end{cauhoids}
